\section{\methodname: An Effective RL Recipe for Code Search}

This section presents the detailed methodology used to train \methodname, covering training data curation and environment construction (\S\ref{sec:data_creation}), the agent scaffold and tools (\S\ref{sec:agent_scaffold}), reward design (\S\ref{sec:reward_design}), and the overall RL training setup (\S\ref{sec:train_algo}).

% We demonstrate an accessible method to produce a competitive code-search agent using open-source artifacts. We train a 14b and a 4b model with this method.

\subsection{Data and Environment Curation}
\label{sec:data_creation}

We describe our approach for curating data and constructing RL environments for training \methodname.

Given a GitHub issue $\mathcal{I}$ in a Python code repository $\mathcal{R}$, we process the ground-truth patch $\mathcal{P}$ which fixes the issue to extract the ground-truth localization target $y^{*}$ at three levels of granularity: edited files, modules, and functions. The LLM agent is tasked with predicting $y^{*}$ given the issue description and the pre-PR repository. 
We build on patch-processing scripts from LocAgent~\citep{chen-etal-2025-locagent} and extend them to handle the following edit operations: detecting additions of member functions and class attributes at the module and file levels; capturing modifications to import statements and global variables at the file level; and ignoring edits to docstrings within functions and classes.\sonicomment{Can we just mention we adapt LocAgent logic and push (non-interesting) details to appendix?}

We curate training instances by processing GitHub issues collected by prior work~\citep{yang2025swesmithscalingdatasoftware, pan2025trainingsoftwareengineeringagents}, originally intended for training agents to fix issues. Note that we discard issues whose PRs create or delete files, as the agent cannot predict the exact filename of newly created files, and we cannot determine the ground-truth modules or functions for deleted files. We further ignore non-Python files (e.g., \texttt{README.md}) modified by the PR from the ground truth as function- and module-level information cannot be extracted from them. Finally, we remove instances with empty issue descriptions.

We construct the RL environment for our agent by simply cloning the pre-PR commit of the repository into a fixed path known to the agent through its prompt. As the localization task does not require executing repository code, we do not require installing project dependencies or use sandboxing/containerization. Notably, since our agent scaffold only requires a bash terminal, the overhead of setting up the environment is \emph{significantly} lower compared to prior approaches whose agents require tools that leverage code dependency graphs, vector databases, module call graphs, or dependency parsers. \sonicomment{TODO: For each complex tool, cite the paper that used this tool}

\subsection{OpenHands-Bash: Our Agent Scaffold}\label{sec:agent_scaffold}

We utilize the OpenHands Software Agent SDK~\cite{wang2025openhandssoftwareagentsdk} to implement the agent scaffold for \methodname given its strong performance on a broad range of software engineering benchmarks, its robust implementation of essential tools (e.g., the terminal), in-built support for parallel tool-calling, and its modular design that facilitates seamless integration with our training backend. 

The agent is primarily equipped with the \texttt{Terminal} tool that supports execution of standard Unix commands like \texttt{rg} (ripgrep), \texttt{find}, \texttt{ls}, \texttt{grep}, \texttt{sed}, etc. To accelerate grep-style code search, we install ripgrep~\footnote{\url{https://github.com/BurntSushi/ripgrep}} in the agent's environment and explicitly prompt the agent to use ripgrep. In addition, the agent is given a \texttt{LocalizationFinish} tool to submit its predicted files, modules, and functions. We refer to this agent scaffold as \textbf{OpenHands-Bash}.

Our initial experiments used the string-based output format from LocAgent~\citep{chen-etal-2025-locagent}, but found it to be sensitive to noisy reward signals due to brittle format validation and parsing logic. We address this by requiring the agent to terminate via the \texttt{LocalizationFinish} tool, which enforces a structured output schema and enables precise parsing and validation, thereby improving the fidelity of the reward signal. Appendix~\ref{appendix:prompts} includes detailed system and user prompts for OpenHands-Bash.

\subsection{Reward Design}\label{sec:reward_design}

For each multi-turn trajectory $\tau$ sampled from the LLM agent, we extract the predicted localization output $y$ from the \texttt{LocalizationFinish} tool, which specifies the predicted files, modules, and functions. Given the ground-truth localization target $y^{\star}$, we compute F1 scores between $y$ and $y^{\star}$ at three granularities, file, module, and function, denoted by $r^{\mathsf{F1-file}}$, $r^{\mathsf{F1-module}}$, and $r^{\mathsf{F1-func}}$, respectively. The reward function is defined as the sum of these three components:
\begin{equation}
\begin{aligned}
    r(\tau, y, y^{\star}) &= r^{\mathsf{F1-file}}(y, y^{\star}) + r^{\mathsf{F1-module}}(y, y^{\star}) \\
    &\quad + r^{\mathsf{F1-func}}(y, y^{\star})
\end{aligned}
\label{eqn:reward}
\end{equation}

During early experiments with \methodname-14B, we observed training collapse characterized by near-zero rewards in later stages. Our analysis revealed that the agent frequently exhausted the maximum number of steps without submitting predictions. To mitigate this issue, we introduce an auxiliary binary reward $r^{\mathsf{turn}}(\tau, k)$ that assigns a reward of 1 if and only if the agent terminates its execution in exactly $k$ turns, where $k$ denotes the step limit. This simple mechanism results in consistent prediction generation within the allotted step limit. Therefore, the final reward function for \methodname-14B is obtained by adding this binary reward to $r(\tau, y, y^{\star})$ in Equation~\ref{eqn:reward}.

We also explored additional auxiliary rewards to incentivize parallel tool calling motivated by SWE-grep~\citep{cognition2025swegrep}, but found these interventions to harm overall performance. Thus, we simply fallback to explicitly prompting the agent to use parallel tool-calling (Appendix \ref{appendix:prompts}).

\subsection{RL Training Algorithm}\label{sec:train_algo}
We implement our training backend using SkyRL~\cite{griggs2025skrylv01}, a modular framework for reinforcement learning with LLMs. It supports asynchronous training~\citep{fu2025areallargescaleasynchronousreinforcement} which improves GPU utilization by parallelizing rollout generation and weight optimization. Specifically, parameter updates are triggered once a sufficient number of rollouts are collected, allowing trajectories sampled from slightly stale model checkpoints. Concretely, trajectories used during the $n^{\text{th}}$ training step may be sampled from model checkpoints that are at most $t$ steps stale, where $t$ denotes the maximum staleness (set to 4 in our experiments). After each optimization step, the updated weights are synchronized with the vLLM~\citep{kwon2023efficient} inference engines, and any in-progress inference requests are terminated to ensure consistency.

We train \methodname models using the Group Sequence Policy Optimization (GSPO) algorithm~\citep{zheng2025groupsequencepolicyoptimization}. \sonicomment{TODO: define all the symbols used below}. The loss is expressed as:
\begin{align}
    \mathcal{J}_{\text{GSPO}}(\theta) &= \nonumber \\  \mathbb{E}_i &\left[ \frac{1}{G} \sum_{i=1}^{G} \min\left( s_i(\theta) \hat{A}_i,\, \text{clip}\left( s_i(\theta), \varepsilon \right) \hat{A}_i \right) \right]
    % \\
    % \nonumber \\
    % \hat{A}_i &= r_i - \text{mean}(\mathbf{r}), \mathbf{r} = \{r_0, ..., r_N\}
    % \label{eq:gspo}
\end{align}
where $s_i$ is the importance ratio derived from sequence likelihood~\cite{zheng-etal-2023-click}:
\begin{align}
    s_i(\theta) = \exp\left( \frac{1}{|y_i|} \sum_{t=1}^{|y_i|} \log \frac{\pi_\theta(y_{i,t} | x, y_{i,<t})}{\pi_{\theta_{\text{old}}}(y_{i,t} | x, y_{i,<t})} \right)
    \label{eq:is}
\end{align}
Following DR. GRPO~\cite{liu2025understandingr1zeroliketrainingcritical}, we remove KL regularization term from the loss function and the standard deviation from the advantage calculation.
\begin{align}
    \hat{A}_i = r_i - \text{mean}(\mathbf{r}), \mathbf{r} = \{r_0, ..., r_N\}
\end{align}
In preliminary experiments, we found GRPO~\cite{shao2024deepseekmathpushinglimitsmathematical} is prone to a collapse later in training, where the tool-calling performance of the model would significantly degrade and would cause zero rewards for most of the rollouts at each training step. \sonicomment{Include more details: we mask loss from trajectories that exhaust max steps without calling the finish tool, we disable KL penalty (cite DAPO), disable entropy loss in optimization}